\section{Experiments}
\label{sec:experiments}

\subsection{Benchmark protocol and baselines}

A useful benchmark must be reproducible, difficult enough to separate methods,
and rich enough to explain the separation. Representative policies therefore
operate on the same garments and initial states under the same closed-loop
budget. Success measures task completion, while distance and difficulty strata
show whether two policies fail in the same way.

\subsection{Does NOCS produce a better trajectory agent?}

NOCS is justified only if it improves the ability to produce successful
physical trajectories, rather than merely making training loss smaller. NOCS
and ordinary-appearance agents therefore share demonstrations, architecture,
optimization, and evaluation cases. Their success and distance across data
scales and garments measure the value of the privileged representation.

\subsection{What does generated data add?}

Generated data can add either new behavior or new appearance. Re-rendering a
fixed trajectory tests visual diversity alone; generating new trajectories
tests state--action diversity; combining both tests their interaction. Varying
garment breadth separately from trajectories per garment then reveals whether
downstream learning benefits from more structural coverage or denser practice
on familiar structures.

\subsection{Where do current policies fail?}

Aggregate success mixes several causes of failure. Conditioning results on
initial distance and garment morphology exposes geometric difficulty. Holding
the physical case fixed while changing texture, background, or lighting exposes
visual sensitivity. Changing validated physical parameters exposes dynamics
sensitivity. Together these views turn the task into a diagnostic benchmark
rather than a single leaderboard number.

\subsection{Do memory and action models help?}

Flattening creates reasons for temporal reasoning that are absent from a
single-frame task. Earlier views can reveal garment identity before it becomes
occluded, and the result of a failed dynamic action indicates how the next
action should change. History, explicit memory, action chunking, RTC, and
world-action models represent different ways to use this information. Their
value is determined by where they improve recovery and difficult cases, not by
their average score alone.

\subsection{Dynamics and real-world transfer}

Dynamic actions couple cloth properties with robot execution. Changes in
friction, stiffness, damping, tracking bandwidth, or action latency can alter
the result of an otherwise identical command. Robustness experiments therefore
test whether a policy observes and corrects physical outcomes instead of
replaying a fixed motion. Real-data experiments hold human data fixed so that
any gain can be attributed to the additional simulated trajectories.
